\paragrafos{Κέντρο συμμετρίας σχήματος}{Β.2.5, «Κέντρο συμμετρίας»}

\begin{anagke}
Η Ρήνη βρίσκει στο συρτάρι μια τράπουλα. «Κοίτα το φύλλο με το καρό: αν το γυρίσω ανάποδα, φαίνεται ακριβώς το ίδιο!». Ο Βασίλης δοκιμάζει με το γράμμα Ν της λέξης «Νίκος»: το γυρίζει μισή στροφή και πάλι Ν. Με το γράμμα Γ όμως δεν συμβαίνει το ίδιο.

Η Αγγελική ρωτάει: «Υπάρχει ένα σημείο, γύρω από το οποίο όταν γυρίζουμε το σχήμα μισή στροφή, αυτό πέφτει πάνω στον εαυτό του;». Η Πολυξένη απαντά: «Όταν υπάρχει, το λέμε κέντρο συμμετρίας του σχήματος».

Σε αυτή την παράγραφο ψάχνουμε ποια σχήματα έχουν κέντρο συμμετρίας και πού βρίσκεται.
\end{anagke}

\begin{lexi}{κέντρο συμμετρίας σχήματος}{κέντρο, συμμετρία}
Ένα σχήμα έχει \lexiword{κέντρο συμμετρίας} αν υπάρχει σημείο $\mathrm{Ο}$ τέτοιο ώστε, όταν το περιστρέψουμε κατά $180^\circ$ γύρω από αυτό, το σχήμα να συμπέσει με τον εαυτό του.
\end{lexi}

\begin{metafrasi}
Μισή στροφή και το σχήμα φαίνεται το ίδιο. & το σχήμα έχει κέντρο συμμετρίας\\
Το συμμετρικό του σχήματος ως προς $\mathrm{Ο}$ είναι το ίδιο το σχήμα. & το $\mathrm{Ο}$ είναι κέντρο συμμετρίας\\
Κάθε σημείο έχει «ταίρι» απέναντι. & για κάθε $\mathrm{Α}$ υπάρχει $\mathrm{Α}'$ στο σχήμα με μέσο το $\mathrm{Ο}$\\
\end{metafrasi}

\tmima{Η θεωρία}

\begin{orismos}{κέντρο συμμετρίας σχήματος}
\textbf{Κέντρο συμμετρίας} σχήματος είναι ένα σημείο $\mathrm{Ο}$ γύρω από το οποίο, αν περιστραφεί το σχήμα κατά $180^\circ$, συμπίπτει με το αρχικό. Τότε το συμμετρικό του σχήματος ως προς το $\mathrm{Ο}$ είναι το ίδιο το σχήμα.
\end{orismos}

Παραδείγματα: το \textbf{παραλληλόγραμμο} έχει κέντρο συμμετρίας το σημείο τομής των διαγωνίων του, ο \textbf{κύκλος} το κέντρο του, το \textbf{ευθύγραμμο τμήμα} το μέσο του. Το \textbf{ισόπλευρο τρίγωνο} και το \textbf{ισοσκελές} δεν έχουν κέντρο συμμετρίας.

\begin{kanonas}{συμμετρική ευθεία}
Το συμμετρικό μιας ευθείας $\varepsilon$ ως προς σημείο $\mathrm{Ο}$ είναι ευθεία $\varepsilon'$ \textbf{παράλληλη} προς την $\varepsilon$. Αν η $\varepsilon$ περνά από το $\mathrm{Ο}$, τότε $\varepsilon'=\varepsilon$.
\end{kanonas}

\begin{center}\small
\begin{tabular}{lcc}
\toprule
Σχήμα & Άξονες συμμετρίας & Κέντρο συμμετρίας\\
\midrule
Ευθύγραμμο τμήμα & $2$ & ναι (μέσο)\\
Ισοσκελές τρίγωνο & $1$ & όχι\\
Ισόπλευρο τρίγωνο & $3$ & όχι\\
Παραλληλόγραμμο & κανένας & ναι\\
Ορθογώνιο & $2$ & ναι\\
Ρόμβος & $2$ & ναι\\
Τετράγωνο & $4$ & ναι\\
Κύκλος & άπειροι & ναι (κέντρο)\\
\bottomrule
\end{tabular}
\end{center}

\begin{prosoxi}
\begin{itemize}[leftmargin=1.2em,itemsep=0.3ex,topsep=2pt]
\item Άξονας συμμετρίας και κέντρο συμμετρίας είναι \emph{ανεξάρτητα}: υπάρχουν σχήματα με το ένα και χωρίς το άλλο (το ισοσκελές τρίγωνο έχει άξονα αλλά όχι κέντρο, το παραλληλόγραμμο κέντρο αλλά όχι άξονα).
\item Ένα σχήμα έχει το πολύ ένα κέντρο συμμετρίας (εκτός από τα άπειρα σχήματα, όπως η ευθεία).
\item Το κέντρο συμμετρίας δεν χρειάζεται να ανήκει στο σχήμα (π.χ. στο δαχτυλίδι).
\end{itemize}
\end{prosoxi}

\tmima{Λυμένα παραδείγματα}

\paradeigma{Το γράμμα}
\tixero Τα γράμματα Ο, Ζ, Μ, Η, Ε.\par
\tithelo Ποια έχουν κέντρο συμμετρίας.\par
\possindeo Κάνω στροφή $180^\circ$ και συγκρίνω: το Ο και το Η συμπίπτουν· το Ζ γυρισμένο φαίνεται Ζ (ναι)· το Μ φαίνεται ανάποδο, ως W (όχι)· το Ε γίνεται «καθρεφτισμένο» (όχι).\par
\apantisi Κέντρο συμμετρίας έχουν τα Ο, Η, Ζ.\par
\elegxos Κάθε σημείο της μισής μορφής έχει ταίρι στην άλλη μισή.

\paradeigma{Το ορθογώνιο}
\tixero Ορθογώνιο $\mathrm{ΑΒΓΔ}$ με διαστάσεις $8$ cm και $6$ cm.\par
\tithelo Το κέντρο συμμετρίας και η απόστασή του από κάθε κορυφή.\par
\possindeo Το κέντρο είναι το σημείο τομής των διαγωνίων. Η διαγώνιος είναι $\sqrt{8^2+6^2}=10$ cm και το κέντρο είναι το μέσο της, άρα απέχει $5$ cm από κάθε κορυφή.\par
\apantisi $\mathrm{ΟΑ}=\mathrm{ΟΒ}=\mathrm{ΟΓ}=\mathrm{ΟΔ}=5$ cm.\par
\elegxos Τα συμμετρικά ως προς $\mathrm{Ο}$ ζεύγη κορυφών ($\mathrm{Α},\mathrm{Γ}$) και ($\mathrm{Β},\mathrm{Δ}$) έχουν το $\mathrm{Ο}$ μέσο.

\paradeigma{Η ευθεία}
\tixero Η ευθεία $\varepsilon$ και σημείο $\mathrm{Ο}$ εκτός αυτής.\par
\tithelo Η συμμετρική ευθεία.\par
\possindeo Παίρνω δύο σημεία $\mathrm{Α},\mathrm{Β}$ της $\varepsilon$, βρίσκω τα $\mathrm{Α}',\mathrm{Β}'$ και φέρνω την ευθεία από αυτά.\par
\apantisi Η $\varepsilon'$ είναι παράλληλη προς την $\varepsilon$.\par
\elegxos Αν το $\mathrm{Ο}$ ήταν πάνω στην $\varepsilon$, η $\varepsilon'$ θα ταυτιζόταν με την $\varepsilon$.

\tmima{Ασκήσεις}

\zoni{Το κατάλαβα}{ποιος έχει κέντρο}

\ex Ποια από τα κεφαλαία γράμματα Α, Η, Ι, Λ, Ν, Ο, Σ, Χ έχουν κέντρο συμμετρίας;

\ex Έχει κέντρο συμμετρίας το ισοσκελές τρίγωνο; Και το ορθογώνιο;

\ex Πόσους άξονες και πόσα κέντρα συμμετρίας έχει ο κύκλος;

\ex Βρες το κέντρο συμμετρίας ενός ρόμβου με διαγώνιους $6$ cm και $8$ cm και πόσο απέχει από κάθε κορυφή του.

\ex Σχεδίασε παραλληλόγραμμο και βρες το κέντρο συμμετρίας του.

\ex Σχεδίασε ευθεία $\varepsilon$ και σημείο $\mathrm{Ο}$ εκτός της. Σχεδίασε τη συμμετρική της.

\ex Σωστό ή λάθος; α) Κάθε σχήμα με άξονα συμμετρίας έχει και κέντρο συμμετρίας. \quad β) Το τετράγωνο έχει κέντρο συμμετρίας. \quad γ) Ο κύκλος έχει άπειρους άξονες συμμετρίας.

\zoni{Το χρησιμοποιώ}{αντικείμενα γύρω μας}

\ex \textbf{Η τράπουλα.} Ο ρόμβος που είναι τυπωμένος στο καρό φύλλο, αν γυρίσεις το φύλλο ανάποδα, μοιάζει ίδιος. Τι σημαίνει αυτό για το κέντρο του ρόμβου;

\ex \textbf{Το πλακάκι.} Ένα ορθογώνιο πλακάκι $30\times20$ cm έχει κέντρο συμμετρίας. Πού βρίσκεται και πόσο απέχει από τις πλευρές του;

\ex \textbf{Η ρόδα.} Μια ρόδα με $8$ ίσα ακτινωτά κομμάτια έχει κέντρο συμμετρίας το κέντρο της. Πόσα ζεύγη «ταίρι» κομμάτια υπάρχουν;

\ex \textbf{Το σήμα κυκλοφορίας.} Το σήμα «STOP» και το «Απαγορεύεται η είσοδος» έχουν κέντρο συμμετρίας; Εξήγησε.

\ex \textbf{Το ταβλί.} Ένα ταμπλό ταβλιού έχει δύο όμοια μισά. Σχεδίασε το κέντρο συμμετρίας.

\zoni{Το σκέφτομαι}{συνδυασμοί συμμετριών}

\ex Βρες σχήμα με κέντρο αλλά χωρίς άξονα συμμετρίας.

\ex Βρες σχήμα με άξονα αλλά χωρίς κέντρο συμμετρίας.

\ex Αν ένα σχήμα έχει δύο κάθετους άξονες συμμετρίας, έχει κέντρο συμμετρίας; Πού;

\ex Εξήγησε γιατί το συμμετρικό ευθείας ως προς σημείο εκτός της είναι παράλληλο. (Βοήθεια: κάθετο τμήμα $\mathrm{ΟΑ}$ από το $\mathrm{Ο}$ προς την $\varepsilon$.)

\zoni{Για το σπίτι}{κέντρα παντού}

\ex Βρες πέντε αντικείμενα στο σπίτι με κέντρο συμμετρίας.

\ex Φτιάξε ένα σχήμα με κέντρο συμμετρίας από διπλωμένο χαρτί.
